\begin{textsample}{NEG Dim 2 – 2009\_c6 – Score -2.00 – 2009\_c6/enade\_2009\_jorn\_09ea6e96cd.txt}  \label{ex:f2_neg_138}
Capítulo 3 Análise Técnica da Prova Este capítulo tem por objetivo apresentar o desempenho dos estudantes inscritos e presentes na prova de COMUNICAÇÃO SOCIAL – JORNALISMO no ENADE/2009.
Para isso, foram calculadas as estatísticas básicas da prova como um todo, bem como as estatísticas das questões relacionadas à Formação Geral e ao Componente Específico.
Nas Tabelas, são apresentadas as seguintes estatísticas : total da população e de presentes, média, erro-padrão da média, desvio-padrão, mediana, nota máxima, nota mínima e coeficiente de assimetria ( ver definição em 1.4.3 ).
Nas estatísticas das questões discursivas, foram detalhadas também as notas nulas : em branco, por protesto ou por insuficiência.
As estatísticas apresentadas neste capítulo contemplam, separadamente, os ingressantes, os concluintes e o total de estudantes, e foram calculadas tendo -se em vista as seguintes agregações : Região Geográfica e País como um todo e ( b ) a Categoria Administrativa e a Organização Acadêmica.
As notas brutas das questões discursivas em cada uma das partes, Formação Geral/Componente Específico, foram calculadas a partir da média simples das notas de cada uma das questões correspondentes.
A primeira questão discursiva da Formação Geral foi desconsiderada e, portanto, a nota bruta referente às questões discursivas de Formação Geral foi igual à nota da questão mantida, a segunda.
Na componente específica a nota bruta correspondeu à média das notas nas três questões.
As notas brutas das questões objetivas de cada uma das partes, Formação Geral/Componente Específico, foram calculadas a partir dos acertos dos estudantes considerando -se uma regra de três para a conversão para a escala de 0 a 100 descontando -se aquelas questões que foram desconsideradas pela comissão ou que não passaram pelo teste do índice de discriminação ponto-bisserial ( ver seção 1.4.6 para a descrição do índice ).
Questões com índices de discriminação inferior a 0,2 foram excluídas do cômputo da nota do estudante.
As notas brutas de Formação Geral e no Componente Específico para cada um dos estudantes foram calculadas como médias ponderadas das suas notas nas questões objetivas e discursivas, ponderação diferente para cada parte : C\_FG = C\_FG^OBJ × 0,6 + C\_FG^DIS × 0,4 C\_CE = C\_CE^OBJ × 0,85 + C\_CE^DIS × 0,15 onde C\_FG é a nota bruta da Formação Geral, C\_OBJ\_FG é a nota bruta nas questões objetivas da Formação Geral, C\_DIS\_FG é a nota bruta nas questões discursivas da Formação Geral, C\_CE é a nota bruta do Componente Específico, C\_OBJ\_CE é a nota bruta nas questões objetivas do Componente Específico, e C\_DIS\_CE é a nota bruta nas questões discursivas do Componente Específico.
A nota bruta da prova, C\_GER, por sua vez, é também uma média ponderada da nota de cada uma das partes, Formação Geral/Componente Específico : C\_GER = 0,25 * C\_FG + 0,75 * C\_CE. 3.1 Estatísticas Básicas da Prova 3.1.1 Estatísticas Básicas Gerais A Tabela 3.1 apresenta as estatísticas básicas da prova por grupo de estudantes.
O total de estudantes inscritos para a prova foi de 24906.
Destes, 75,4 \% compareceram, sendo que a presença foi maior entre concluintes ( 85,5 \% ) do que entre ingressantes ( 66,6 \% ).
A média geral da prova foi 46,6, sendo que os ingressantes obtiveram média mais baixa, 44,2, que os concluintes, cuja média foi 48,6.
O desvio-padrão geral foi de 17,6, sendo o do grupo de ingressantes menor ( 16,9 ), indicando que, em contraposição ao grupo dos concluintes ( 18,0 ), aqueles possuíram uma distribuição mais homogênea nas notas.
A nota máxima foi 93,5, obtida por pelo menos um ingressante ao passo que a maior nota obtida por um concluinte foi 90,1.
Os coeficientes de assimetria para os dois grupos considerados, o de concluintes e o de ingressantes, bem como para a população como um todo foram todos negativos, indicando uma maior concentração à direita da distribuição e um maior espalhamento à esquerda.
Uma comparação entre o desempenho de ingressantes e concluintes pode ser realizada a partir dos dados apresentados no Gráfico 3.1.
Cumpre destacar que as notas dos concluintes, de um modo geral, foram mais elevadas do que as dos ingressantes, cuja distribuição de notas está à esquerda dos primeiros.
A ideia de melhor desempenho é reforçada pelas notas medianas dos dois grupos : 50,8 entre os concluintes e 45,5 entre os ingressantes.
Nas faixas de notas até 50 pontos o percentual de ingressantes é maior que o dos concluintes, situação esta que se inverte para as faixas de notas superiores a 50 pontos.
A exceção é o primeiro intervalo de 0 a 10 que inclui uma quantidade maior de concluintes do que de ingressantes.
Constatou -se que aproximadamente 39,3 \% dos ingressantes situam -se nas faixas de notas superiores a 50,0 pontos, por oposição a 52,1 \% dos concluintes.
O Gráfico 3.2 apresenta informações referentes ao desempenho geral de ingressantes e concluintes, comparando os resultados considerando a Região, Categoria Administrativa e Organização Acadêmica.
Levando -se em conta as notas médias dos estudantes em cada Região, observa -se que, há nos dois grupos, uma variação expressiva entre as Regiões.
As notas médias mais elevadas entre os concluintes ocorreram na Região Sul ( 55,05 ), acima da média nacional ( 48,62 ).
Em relação aos ingressantes, as médias mais altas foram obtidas, também nas Regiões Sul ( 48,65 ) e Nordeste ( 45,35 ) nas demais Regiões a média foi inferior à média nacional ( 44,23 ).
Na Região Norte foi obtida a média de notas mais baixa entre os concluintes ( 42,16 ) e os ingressantes ( 38,29 ).
Cursos nesta área só foram oferecidos em instituições federais, estaduais, municipais e particulares.
Levando -se em conta os agrupamentos dos estudantes em categorias administrativas, observa -se que a pontuação média mais elevada entre os concluintes foi encontrada em instituições de origem municipal ( 52,08 ) – acima da média nacional, se observa que a pontuação média mais elevada entre os ingressantes foi encontrada em instituições de origem federal ( 48,33 ) – também acima da média nacional.
Tendo como foco as notas médias dos estudantes agrupados por Organização Acadêmica, notamos que universidades apresentam as maiores médias ( concluintes – 49,08 e ingressantes – 45,32 ), ao passo que as piores notas médias foram observadas para concluintes nas faculdades e para ingressantes nos centros universitários. 3.1.2 Estatísticas Básicas em Formação Geral A Tabela 3.2 apresenta as estatísticas básicas de Formação Geral por grupo de estudantes.
O total de estudantes inscritos para a prova foi de 24906.
Destes, 75,4 \% compareceram, sendo que a presença foi maior entre concluintes ( 85,5 \% ) do que entre ingressantes ( 66,6 \% ).
A média geral da parte da prova de Formação Geral foi de 43,12, sendo que os ingressantes obtiveram média mais baixa, 42,01, que os concluintes, cuja média foi 44,10.
O desvio-padrão geral foi de 22,3, sendo o do grupo de ingressantes menor ( 21,9 ), do que o do grupo dos concluintes ( 22,5 ), aqueles possuíram uma distribuição mais homogênea nas notas.
A nota máxima foi 91,4, obtida por estudantes de ambos os grupos.
Para esta parte, os coeficientes de assimetria são positivos para os ingressantes e nulos para os demais.
O Gráfico 3.3 permite a comparação entre concluintes e ingressantes relativamente ao desempenho na parte da prova que aborda a Formação Geral.
Nesta parte, se excluindo as notas nulas, tal como na prova como um todo, também os concluintes apresentaram desempenho superior, embora seja perceptível um maior equilíbrio.
A mediana dos concluintes foi igual à dos ingressantes, já a média dos concluintes é superior à dos ingressantes.
Dos ingressantes, 36,2 \% alcançaram nota superior a 50 ; dentre os concluintes, o percentual dos estudantes nesta faixa de notas é 40,4 \%.
No Gráfico 3.4 são apresentadas as informações referentes ao desempenho dos ingressantes e dos concluintes na parte de Formação Geral, em diferentes agregações como Região, Categoria Administrativa e Organização Acadêmica.
Considerando -se as notas médias dos estudantes segundo cada Região, observou -se que, tanto no grupo dos ingressantes quanto no grupo dos concluintes, há diferenças marcantes.
As médias variam, entre os concluintes, de 33,51 na Região Norte a 53,01 na Região Sul.
Entre os ingressantes, os valores foram de 31,60 na Região Norte a 48,48 na Região Sul.
A maior dispersão de notas dos estudantes de uma dada Região foi encontrada no Norte tanto para os concluintes quanto para os ingressantes.
Em todas as Regiões os concluintes apresentaram médias maiores do que os ingressantes.
Todas as Regiões apresentaram concluintes.
Tendo como foco as categorias administrativas, destacou -se que todas as instituições estavam representadas.
Sendo que as notas médias das federais foram superiores para ingressantes ( 47,69 ), sendo inclusive superiores às notas dos seus concluintes.
Por outro lado, as notas médias das municipais foram superiores para concluintes ( 48,32 ).
Os estudantes das instituições particulares se apresentam mais homogêneos, já que o desvio padrão para concluintes e ingressantes foi o menor dos valores encontrados.
Considerando -se o tipo de Organização Acadêmica, somente as Universidades apresentaram os maiores valores, tanto para concluintes ( 44,72 ) quanto para ingressantes ( 43,47 ).
Os estudantes dos Centros Universitários se apresentaram menos homogêneos, já que o desvio padrão é o maior encontrado em ambos os grupos.
Os Centros Federais de Educação Tecnológica e os Institutos Federais de Educação, Ciência e Tecnologia não estavam representados nesta área. 3.1.3 Estatísticas Básicas em Componente Específico A Tabela 3.3 apresenta as estatísticas básicas referentes ao Componente Específico da prova, tendo sua média geral, de 47,8, superior ao desempenho da prova de COMUNICAÇÃO SOCIAL – JORNALISMO como um todo cuja média foi 46,6, como apresentado na Tabela 3.1, que por sua vez foi superior ao desempenho na Formação Geral, que teve média igual a 43,12, conforme mostrado na Tabela 3.2.
Os concluintes apresentaram um desempenho médio ( 50,19 ) superior ao dos ingressantes ( 45,02 ), havendo ainda indícios de maior homogeneidade entre as notas médias dos ingressantes do que entre as dos concluintes, como sugeriu o desvio padrão, menor no grupo dos ingressantes do que no dos concluintes.
As notas máximas obtidas por concluintes ( 96,2 ) foram inferiores às obtidas por ingressantes ( 97,1 ).
As distribuições de notas para cada grupo se apresentaram com uma inclinação para notas maiores como evidenciado pelos coeficientes de assimetria negativos.
A situação de assimetria é confirmada pelo Gráfico 3.5 que proporciona uma comparação do desempenho entre concluintes e ingressantes em relação ao Componente Específico da prova.
Claramente a distribuição das notas médias dos ingressantes, à direita daquela dos concluintes, evidenciou o percentual mais elevado de notas médias nas faixas mais baixas para este grupo, com exceção do primeiro, que como já comentado, incluiu uma maior proporção de zeros entre os concluintes.
As medianas das distribuições das notas dos ingressantes e dos concluintes foram 46,6 e 52,6, respectivamente, \textbf{enfatizando} o melhor desempenho dos concluintes nesta parte da prova.
Assim como os Gráficos 3.2 e 3.4, o Gráfico 3.6 apresenta uma comparação dos resultados em relação às Regiões, Categoria Administrativa e Organização Acadêmica, agora levando em conta o desempenho dos ingressantes e dos concluintes no Componente Específico da prova.
Tendo como foco as notas médias dos estudantes, em determinada Região, observa -se que, entre os concluintes, as mais \textbf{elevadas} foram encontradas na Região Sul ( 55,78 ), valores acima da média nacional para esta parte ( 50,18 ).
Entre os ingressantes, as maiores notas ocorreram nas Regiões Nordeste ( 46,10 ) e Sul ( 48,75 ), todos acima da média nacional ( 45,02 ).
Novamente, para a análise segundo a Categoria Administrativa, todas as entidades foram representadas.
Observa -se comportamento semelhante àquele da parte relativa da Formação Geral, ou seja, as médias mais elevadas entre os concluintes foram encontradas em instituições de origem municipal ( 53,39 ) e em instituições de origem federal ( 48,59 ) entre os ingressantes.
Os estudantes das instituições particulares apresentam uma homogeneidade, com desvio padrão das notas iguais para os concluintes e para os ingressantes.
Quanto à Organização Acadêmica, o comentário é o mesmo já feito para a nota bruta e referente à Formação Geral : somente as Universidades, Centros Universitários e Faculdades foram representados neste curso.
Considerando -se o tipo de Organização Acadêmica, entre os concluintes e ingressantes, Universidades apresentaram os maiores valores, seguidas dos Centros Universitários e das Faculdades. 3.2 Análise das Questões Objetivas 3.2.1 Formação Geral A Tabela 3.4 apresenta as estatísticas básicas relativas às oito questões objetivas da Formação Geral dos estudantes.
Adiantamos que a primeira questão foi desconsiderada pelo critério da correlação ponto-bisserial.
Como pode ser observado na Tabela 3.4, os estudantes concluintes obtiveram desempenho superior, com nota média de 52,47, em relação aos estudantes ingressantes cuja nota média foi 50,48.
O desvio padrão entre os concluintes foi 22,7 e entre os ingressantes 22,4, indicaram maior homogeneidade para o grupo dos ingressantes.
A amplitude total das notas foi a mesma para ambos os grupos : nota mínima de 0,00 e máxima de 85,7.
A mediana foi igual para ambos os grupos : 57,1.
A assimetria, medida pelo coeficiente de assimetria ( skewness ) foi um pouco menor em módulo para os ingressantes ( -0,4 ) do que para os concluintes ( -0,6 ), ainda que ambos fossem negativos indicando uma assimetria para a esquerda, fato reforçado pelos valores menores das médias dos dois grupos quando comparados à mediana correspondente.
A Tabela 3.5 apresenta a classificação das questões objetivas da Formação Geral, segundo o índice de facilidade ( questões em vermelho foram desconsideradas pelo critério ponto-bisserial ).
Como as questões foram comuns a todas as áreas, a tabela apresenta a classificação considerando -se os concluintes de todas as áreas e aqueles do curso de COMUNICAÇÃO SOCIAL – JORNALISMO.
De acordo com os índices obtidos, as questões objetivas da prova foram assim avaliadas para a área : das oito questões, três tiveram entre 60 \%, inclusive, e 85 \% de acertos e foram classificadas como fáceis ( somente uma para o A Tabela 3.6 descreve os resultados das análises das questões objetivas relativas à Formação Geral, segundo o poder de discriminação, utilizando -se para tal o índice de discriminação ponto-bisserial.
Como as questões foram comuns a todas as áreas, a tabela apresenta a classificação considerando -se os concluintes de todas as áreas e aqueles do curso de COMUNICAÇÃO SOCIAL – JORNALISMO.
A desconsideração da questão segundo este critério considerou o universo de todos os concluintes.
Três das oito questões apresentaram índices iguais ou acima de 0,40 ( nenhuma para a população de concluintes como um todo ), sendo, portanto, classificadas como muito boas nesta característica, para este grupo de estudantes, duas foram classificadas como boas com índice entre 0,30, inclusive, e 0,40 ( três para os concluintes como um todo ) e duas foram classificadas como médias, com índice entre 0,20, inclusive, e 0,30 ( quatro para os concluintes como um todo ).
Uma das questões foi classificada como fraca, tanto para a população de concluintes como um todo como para a de concluintes de COMUNICAÇÃO SOCIAL – JORNALISMO e foi desconsiderada pelo critério ponto-bisserial. total de concluintes ) ; outras duas foram consideradas medianas, alcançando entre 40 \%, inclusive, e 60 \% de acertos ( três para o total de concluintes ) ; duas obtiveram entre 15 \%, inclusive, e 40 \% de acertos, sendo consideradas difíceis ( três para o total de concluintes ).
Por fim, apenas uma foi classificada como muito difícil tento obtido percentual menor do que 15 \% de acertos ( uma para o total de concluintes ), coincidentemente a questão desconsiderada pelo critério ponto-bisserial.
Nenhuma questão foi classificada na categoria muito fácil.
A Tabela 3.6 descreve os resultados das análises das questões objetivas relativas à Formação Geral, segundo o poder de discriminação, utilizando -se para tal o índice de discriminação ponto-bisserial.
Como as questões foram comuns a todas as áreas, a tabela apresenta a classificação considerando -se os concluintes de todas as áreas e aqueles do curso de COMUNICAÇÃO SOCIAL – JORNALISMO.
A desconsideração da questão segundo este critério considerou o universo de todos os concluintes.
Três das oito questões apresentaram índices iguais ou acima de 0,40 ( nenhuma para a população de concluintes como um todo ), sendo, portanto, classificadas como muito boas nesta característica, para este grupo de estudantes, duas foram classificadas como boas com índice entre 0,30, inclusive, e 0,40 ( três para os concluintes como um todo ) e duas foram classificadas como médias, com índice entre 0,20, inclusive, e 0,30 ( quatro para os concluintes como um todo ).
Uma das questões foi classificada como fraca, tanto para a população de concluintes como um todo como para a de concluintes de COMUNICAÇÃO SOCIAL – JORNALISMO e foi desconsiderada pelo critério ponto-bisserial.
Na sequência, no Quadro 3.1, estão especificadas, por questão, as habilidades aferidas pelas questões de Formação Geral da prova da COMUNICAÇÃO SOCIAL – JORNALISMO, bem como os resultados da análise relativa aos índices de facilidade e de discriminação de cada questão.
O índice de facilidade para a área de JORNALISMO teve uma variação considerável, oscilando entre 0,09 até 0,82 e o de discriminação ficou entre 0,05 e 0,57.
A questão 8 aferiu as habilidades de “ Ler e interpretar textos ”, “ Analisar e criticar informações ”, “ Extrair conclusões por indução e/ou dedução ”, “ Estabelecer relações, comparações e contrastes em diferentes situações ”, “ Detectar contradições ”, “ Fazer escolhas valorativas avaliando consequências ”, “ Questionar a realidade ”, e “ Argumentar coerentemente ” foi a que apresentou, para a população de concluintes de COMUNICAÇÃO SOCIAL – JORNALISMO, maior poder discriminatório, ambos com índice de 0,57.
A questão mais fácil, com 82 \% de acertos, foi também a de número 8.
De outra parte, a questão com menor índice de facilidade e de discriminação em Formação Geral foi a questão 1, que foi desconsiderada, e que deveria aferir as habilidades de “ Analisar informações ” ; “ Extrair conclusões ” e “ Propor soluções para situações-problema ”.
O Gráfico 3.7 mostra, a título de ilustração, o comportamento das respostas à questão 7 de Formação Geral para os estudantes concluintes da área de JORNALISMO como função da nota geral na prova.
Trata -se de uma questão considerada fácil para os concluintes de COMUNICAÇÃO SOCIAL – JORNALISMO e, de dificuldade média para o conjunto de concluintes de todas as áreas.
O poder de discriminação foi muito bom para os estudantes da área e bom para o conjunto de todos os concluintes.
Neste Gráfico, cada uma das cinco curvas representou o percentual de respostas em determinada alternativa da questão, em função do escore total dos estudantes concluintes da área de JORNALISMO.
A curva azul marinho corresponde à alternativa A, a correta para esta questão.
Assim, observou -se que entre os estudantes com notas mais baixas – desempenho geral mais fraco na prova como um todo –, houve a escolha de alternativas incorretas, especialmente a alternativa B – curva verde claro seguida da alternativa C – curva em vermelho, além da alternativa correta.
A proporção de estudantes que selecionaram a resposta correta A, aumentou com o desempenho até o grupo de estudantes com notas no intervalo ( 75,0 ; 87,5 ), chegando ao acerto de quase 100 \%, tal percentual decaiu para os estudantes com notas pelo menos 87,5.
Esta análise permitiu verificar como a questão discriminou entre os grupos de desempenho.
Os Gráficos relativos às demais questões de Formação Geral constam do Anexo I. 3.2.2 Componente Específico A Tabela 3.7 apresenta as estatísticas básicas em relação às questões objetivas de Componente Específico.
Como pode ser observado, os estudantes concluintes obtiveram desempenho superior, com uma média de 53,45 pontos, em relação aos estudantes ingressantes com uma média de 48,04 pontos.
Os desvios-padrão indicaram que a variabilidade entre os concluintes foi maior ( 19,7 ) do que entre os ingressantes ( 18,5 ).
Os coeficientes de assimetria para os dois grupos foram negativos indicando uma inclinação para a direita.
A assimetria foi, porém, maior entre os concluintes.
A análise das questões objetivas de Componente Específico em COMUNICAÇÃO SOCIAL – JORNALISMO, em relação ao grau de facilidade da prova ( ver Tabela 3.8 ), mostrou que, das 27 questões, 8 encontraram -se na categoria difícil ( entre 15 \% e 40 \%, não incluso, de acertos ), 9 foram de nível médio ( entre 40 \% e 60 \%, não incluso, de acertos ), cinco estão na categoria fácil ( entre 60 \% e 85 \%, não inclusos, de acertos ) e nenhuma teve a classificação muito fácil ( a partir de 85 \% de acertos ).
Nenhuma questão teve percentual de acertos menor do que 15 \% ( classificação muito difícil ) e cinco questões foram desconsideradas pela Comissão Assessora da Área de JORNALISMO.
Das 27 questões objetivas relativas ao Componente Específico ( ver Tabela 3.9 ), onze obtiveram o índice de discriminação muito bom, quatro obtiveram o índice de discriminação bom e 6 obtiveram nível de discriminação médio.
O índice fraco foi alcançado por uma questão, eliminando -a.
Cinco questões foram desconsideradas pela comissão.
Na sequência, no Quadro 3.2 estão especificadas, por questão, as habilidades aferidas pelas questões de Componente Específico da prova, bem como os resultados da análise relativa aos índices de facilidade e de discriminação de cada questão.
Das 27 questões objetivas relativas ao Componente Específico ( ver Tabela 3.9 ), onze obtiveram o índice de discriminação muito bom, quatro obtiveram o índice de discriminação bom e 6 obtiveram nível de discriminação médio.
O índice fraco foi alcançado por uma questão, eliminando -a.
Cinco questões foram desconsideradas pela comissão.
Na sequência, no Quadro 3.2 estão especificadas, por questão, as habilidades aferidas pelas questões de Componente Específico da prova, bem como os resultados da análise relativa aos índices de facilidade e de discriminação de cada questão.
O Gráfico 3.8 mostra, a título de ilustração, o comportamento da questão 11 do Componente Específico para os estudantes concluintes da área de JORNALISMO como função da nota geral na prova5.
Trata -se de uma questão considerada fácil para os concluintes de COMUNICAÇÃO SOCIAL – JORNALISMO.
O poder de discriminação foi muito bom para os estudantes da área.
Neste Gráfico, cada uma das cinco curvas representa o percentual de respostas em determinada alternativa da questão, em função da nota dos estudantes concluintes de JORNALISMO.
A curva verde claro correspondeu à alternativa B, a correta para esta questão.
Assim, observou -se que entre os estudantes com escores mais baixos – desempenho geral mais fraco na prova como um todo –, houve a escolha de alternativas incorretas, entre elas a alternativa E – curva em roxo e a alternativa A – curva em azul marinho e a alternativa C – curva em vermelho.
A proporção de estudantes que selecionaram a resposta correta B aumentou com o desempenho deste grupo de estudantes chegando ao acerto de 100 \% dos estudantes com notas pelo menos 87,5.
Esta análise permitiu verificar como a questão discriminou entre os grupos de desempenho.
Os Gráficos relativos às demais questões do Componente Específico constam do Anexo I. 3.3 Análise das Questões Discursivas 3.3.1 Formação Geral A Tabela 3.10 apresenta as estatísticas básicas da prova por grupo de estudantes.
A análise dos resultados de desempenho dos estudantes nas questões discursivas que são apresentados na tabela 3.10 e no gráfico 3.9, mostrou que as notas foram mais baixas no conjunto dessas questões que no das questões objetivas.
O total de estudantes inscritos para a prova foi de 24906.
Destes, 75,4 \% compareceram, sendo que a presença foi maior entre concluintes ( 85,5 \% ) do que entre ingressantes ( 66,6 \% ).
A média geral da prova foi de 30,60, sendo que os ingressantes obtiveram média mais baixa, de 29,42, que os concluintes, cuja média foi 31,64.
O desvio, padrão geral foi de 36,5, sendo o do grupo de ingressantes menor ( 35,9 ), em relação ao grupo dos concluintes ( 36,9 ), aqueles possuíram uma distribuição mais homogênea nas notas.
A nota máxima foi 100,0, atingida por concluintes e ingressantes.
Os coeficientes de assimetria para os dois grupos considerados, o de concluintes e o de ingressantes, bem como para a população como um todo foram todos positivos, indicando uma maior concentração à esquerda da distribuição.
Para todos os grupos mais de metade dos estudantes presentes obtiveram nota zero nas questões discursivas da Formação Geral.
Das 10431 provas com nota zero, 540 ( 5,2 \% ) continham um protesto, 3155 ( 30,2 \% ) estavam em branco e 6736 ( 64,6 \% ) foram respondidas, porém de forma inadequada.
Entre os concluintes, a maioria ( 64,3 \% ) dentre aqueles com nota zero apresentaram as provas respondidas de forma inadequada.
Entre os ingressantes, 64,9 \% das notas nulas decorreram de correção.
Uma comparação entre o desempenho de ingressantes e concluintes pode ser realizada a partir dos dados apresentados no Gráfico 3.9.
As notas dos concluintes, de um modo geral, foram ligeiramente mais \textbf{elevadas} do que as dos ingressantes, cuja distribuição de notas está à direita dos primeiros.
A ideia de melhor desempenho foi reforçada pela proporção de notas nulas : 53,8 \% entre os concluintes e 57,5 \% entre os ingressantes.
Na sequência, discorrer -se-á, especificamente, sobre uma das questões discursivas de Formação Geral, a número 10 já que a número 9 foi desconsiderada.
Em especial, serão analisados os conteúdos englobados na questão e os resultados obtidos pelos estudantes.
Ao final, são apresentados os comentários da banca corretora a respeito do observado nas respostas dos estudantes, suas impressões e conclusões.
Cumpre esclarecer que, tendo em vista que as questões discursivas de Formação Geral foram padronizadas, ou seja, constam de todas as provas, os comentários da banca foram os mesmos para todas as áreas acadêmicas, eis que direcionados a todos os estudantes que participaram do ENADE/2009.
A seguir, serão analisados os desempenhos de ingressantes e concluintes da área de COMUNICAÇÃO SOCIAL – JORNALISMO na questão discursiva de Formação Geral, não desconsidera, do ENADE/2009, comparando os resultados obtidos com a habilidade exigida na questão. 3.3.1.1 Análise da Questão Discursiva 9 de Formação Geral DESCONSIDERADA 3.3.1.2 Análise da Questão Discursiva 10 de Formação Geral A Questão Discursiva nº. 10 buscou avaliar as habilidades : • Ler e interpretar textos ; • Analisar e criticar informações ; • Extrair conclusões por indução e/ou dedução ; • Estabelecer relações, comparações e contrastes em diferentes situações ; • Fazer escolhas valorativas avaliando consequências ; • Questionar a realidade ; • Argumentar coerentemente.
O Quadro 3.3 apresenta as habilidades/competências, temas predominantes e estatísticas básicas da questão 10 de Formação Geral por grupo de estudantes.
Cumpre notar que como a questão 9 foi desconsiderada, a nota referente às questões discursivas da Formação Geral é idêntica à nota da questão 10.
A média geral da prova foi de 30,60, sendo que os ingressantes obtiveram média mais baixa, de 29,42, que os concluintes, cuja média foi 31,64.
O desvio, padrão geral foi de 36,5, sendo o do grupo de ingressantes menor ( 35,9 ), em relação ao grupo dos concluintes ( 36,9 ), aqueles possuíram uma distribuição mais homogênea nas notas.
A nota máxima foi 100,0, atingida por concluintes e ingressantes.
Os coeficientes de assimetria para os dois grupos considerados, o de concluintes e o de ingressantes, bem como para a população como um todo foram todos positivos, indicando uma maior concentração à esquerda da distribuição.
Para todos os grupos, mais de metade dos estudantes presentes tiveram nota zero nesta questão discursiva de Formação Geral.
Das 10431 provas com nota zero, 540 ( 5,2 \% ) continham um protesto, 3155 ( 30,2 \% ) estavam em branco e 6736 ( 64,6 \% ) foram respondidas, porém de forma inadequada.
Entre os concluintes, a maioria ( 64,3 \% ) dentre aqueles com nota zero estava com a prova respondida de forma inadequada.
Entre os ingressantes, 64,9 \% das notas nulas decorreram do desempenho do próprio estudante verificadas na correção.
Uma comparação entre o desempenho de ingressantes e concluintes para esta questão foi exatamente igual à já realizada para a nota discursiva em Formação Geral a partir dos dados apresentados no Gráfico 3.9. 3.3.1.3 Comentários sobre as respostas à Questão Discursiva 10 Apresentação A questão 10 de Formação Geral apresentou um questionamento capaz de analisar a habilidade do estudante em reconhecer o que é um direito social, optar por um direito social vigente que possa diminuir as desigualdades de renda familiar no Brasil e discutir os argumentos em favor do direito social escolhido como suporte à resposta.
Para análise das respostas apresentadas pelos estudantes, houve uma leitura inicial bastante criteriosa.
Foram considerados para efeito de pontuação, os seguintes argumentos : • Acesso à educação pública, gratuita e de qualidade, garantindo ao cidadão melhores condições de vida ; • Assistência à saúde, em seu contexto mais amplo, garantindo os cuidados essenciais de saúde, baseados em técnicas apropriadas ao alcance de todas as pessoas, contando com a participação da população ; • Permanência do estudante na escola, em todos os níveis escolares ( da educação infantil à educação superior, possibilitando a qualificação profissional ), o que garante melhores condições de trabalho e remuneração.
Considerou -se a qualificação e o desempenho docente, desenvolvendo nos estudantes melhores expectativas de aprendizagem e interesse pelo conteúdo da escola.
Os professores desenvolvendo metodologias adequadas e conhecimento pleno quanto a questões de aprendizagem ; • Condições dignas de trabalho, com remuneração que garanta qualidade de vida.
Foram consideradas as condições de trabalho, a preparação profissional para o trabalho e os cursos de aperfeiçoamento, oportunizando momentos de aprendizagem e desenvolvimento, a partir de situações desencadeadoras ; • Ser proprietário do imóvel em que reside, reduzindo os gastos com aluguel, garantindo melhores condições de vida, equilíbrio financeiro e construção do bem, estar familiar.
Desempenho dos estudantes na questão É importante salientar que havia respostas que desenvolveram adequadamente o questionamento apresentado, com fundamentação baseada em argumentos coerentes.
Contudo, em um número considerável de respostas, havia apenas um argumento e não dois, conforme fora solicitado.
Muitos estudantes informaram sobre problemas da educação, de um modo geral.
Algumas observações mencionadas pelos corretores durante o processo de correção foram : • há preocupação com uma educação de qualidade, com sugestões de cursos profissionalizantes e métodos preparatórios de aulas. • a educação é revelada como de baixa qualidade, exaltando -se treinamentos e cursos de atualização para a classe docente. • a saúde como essência para a vida.
O SUS é colocado como parte integrante de direito social, garantindo às pessoas fácil acesso. • oportunidade de emprego como forma de ascensão social, direito, cidadania. • os políticos brasileiros como causadores do despreparo educacional.
As respostas analisadas apresentaram norma culta, coesão e coerência razoáveis, que garantiram a compreensão da mensagem desenvolvida pelos estudantes.
Os erros ortográficos foram poucos e não comprometeram a discussão sobre o assunto e sobre os pontos de vista abordados.
Quanto à coerência, houve um desvio quanto ao objetivo da questão, já que alguns estudantes, através dos seus argumentos retrataram reivindicações, reclamações e situações de problemas enfrentados pelos brasileiros, citando golpes políticos.
Porém, mesmo quando se tratava desta avaliação, os estudantes foram autênticos através das explicações.
Houve erros quanto à concordância verbal e nominal.
Houve também, dificuldade em utilizar os conectores corretamente.
Alguns estudantes, por falta de ideias, apresentaram em suas respostas redundância e prolixidade.
O desempenho foi considerado satisfatório, pois a maioria dos estudantes soube responder ao questionamento.
Houve domínio do conteúdo, pois as ideias, os argumentos e a originalidade corresponderam ao esperado. 3.3.2 Componente Específico Nas questões discursivas do Componente Específico – questões de números 38, 39 e 40 –, a diferença de desempenho entre os grupos não foi muito acentuada, inclusive menor do que para as questões discursivas de Formação Geral.
A Tabela 3.11 contém as estatísticas da média destas questões.
Observa -se que no grupo dos concluintes a nota média alcançada foi 32,0,0 enquanto que no grupo dos ingressantes foi 28,24, ambas menores do que as obtidas nas questões discursivas de Formação Geral, que foram 29,42 para os ingressantes e 31,64 para os concluintes ( Tabela 3.10 ).
As médias foram também menores dos que as obtidas para as questões objetivas do Componente Específico ( Tabela 3.7 ) : respectivamente 48,04 e 53,45 para ingressantes e concluintes.
A análise de cada uma destas questões será feita nas subseções seguintes.
Os valores positivos para a medida de assimetria indicaram uma distribuição com cauda para a direita, ou seja, mais concentrada à esquerda, com menores valores.
Dos 18783 presentes, 23,5 \% obtiveram nota zero.
A tabela 3.12 apresenta um detalhamento da situação destes estudantes.
A maior parte dos estudantes ( 50,2 \% ) não fez nenhuma das 3 questões.
Esta proporção foi menor entre os concluintes ( 50,0 \% ) do que entre os ingressantes ( 50,4 \% ).
A combinação de questões em branco e preenchidas correspondeu a 65,2 \% das provas.
Cerca de 1 \% no total das provas correspondeu a 3 questões com protesto.
O Gráfico 3.10 apresenta a distribuição das notas nas questões discursivas do Componente Específico por grupo de estudantes.
A moda se situou no primeiro intervalo, aquele com valores entre 0 e 10 incluindo os extremos do intervalo.
A impressão de assimetria informada pelos coeficientes de assimetria positivos ( Tabela 3.11 ) foi reforçada no gráfico pela concentração no primeiro intervalo. 3.3.2.1 Análise da Questão Discursiva 38 de Componente Específico A questão 38, cujas habilidades aferidas encontram -se descritas no Quadro 3.4, que ainda contém outras informações, seguido da Tabela 3.13 e do Gráfico 3.11, foi aquela em que os grupos de estudantes observados obtiveram o melhor desempenho entre as questões discursivas do Componente Específico.
A nota média do grupo dos ingressantes foi 33,79, e a dos concluintes 36,04.
O percentual de notas zero obtidas, o mais baixo das questões discursivas do Componente Específico, nos dois grupos foi de 28,2 \%, ligeiramente menor entre os ingressantes ( 28,1 \% ) do que entre os concluintes ( 28,2 \% ).
Mais de 70 \% dos concluintes e dos ingressantes que tiraram zero deixaram a questão em branco ( ver Tabela 3.13 ).
Uma pequena minoria ( 5,0 \% ) utilizou a questão para protestar.
Apesar do baixo desempenho, em ambos os grupos a nota máxima de 100,0 pontos foi alcançada.
Os coeficientes de assimetria positivos indicaram uma inclinação para a esquerda e cauda para a direita.
O Gráfico 3.11 apresenta a distribuição das notas na questão 38 discursiva do Componente Específico por grupo de estudantes.
A moda se situou no primeiro intervalo, aquele com valores entre 0 e 10 incluindo os extremos do intervalo.
A impressão de assimetria informada pelos coeficientes de assimetria positivos ( Quadro 3.4 ) foi reforçada no gráfico pela concentração no primeiro intervalo. 3.3.2.2 Comentários sobre as respostas à Questão Discursiva 38 A questão apresentou -se com nível médio, possibilitando aos estudantes expressarem seus conhecimentos relativos às técnicas de produção e linguagem em jornalismo impresso, aprendidos na faculdade.
Percebe -se claramente a distinção de respostas dos concluintes, pois conseguiram responder à questão de forma parcial ou completa dentro dos preceitos jornalísticos.
Foram solicitadas duas respostas, e a primeira que dizia respeito à comparação entre as linhas editoriais foi mais frequente do que as propostas de soluções inovadoras, tendo sido mais atendida pelos estudantes.
A chave de correção sugere que o professor, corretor considere como aspecto a ser pontuado uma análise crítica das diferentes linhas editoriais dos dois jornais delineados.
No entanto, ser objetivo, cuidadoso na checagem da informação, ouvir quantos lados da questão forem necessários e não deixar, se envolver por pressões externas ao trabalho jornalístico em si são, por exemplo, leis elementares e comuns a qualquer tipo de linha editorial.
Veículo popular ou de elite não define diferença de atuação em termos de qualidade profissional.
Tanto assim que bons repórteres dos melhores veículos passaram por outros jornais com menor estrutura antes, e a categoria é unânime em dizer que trabalhar na Editoria de Polícia, a mais valorizada nos jornais populares, é uma das melhores escolas de reportagem que há. 3.3.2.3 Análise da Questão Discursiva 39 de Componente Específico A questão 39, cujas habilidades aferidas encontram -se descritas no Quadro 3.5, que ainda contém outras informações, seguido da Tabela 3.14 e do Gráfico 3.12, foi aquela em que os grupos de estudantes observados obtiveram o desempenho intermediário entre as questões discursivas do Componente Específico.
A nota média do grupo dos ingressantes foi 29,44, e a dos concluintes 32,98.
O percentual de notas zero obtidas nos dois grupos foi 48,7 \% entre os ingressantes e 45,9 \% entre os concluintes.
Das provas com notas nulas nesta questão, aproximadamente 1/3 foi deixada em branco : 32,9 \% entre os ingressantes e 33,1 \% entre os concluintes ( ver Tabela 3.14 ).
Entre os com nota nula, uma pequena minoria ( 2,0 \% ) utilizou a questão para protestar.
Apesar do baixo desempenho, em ambos os grupos a nota máxima de 100,0 pontos foi alcançada.
Os coeficientes de assimetria positivos indicaram uma inclinação para a esquerda e cauda para a direita.
O Gráfico 3.12 apresenta a distribuição das notas na questão 39 discursiva do Componente Específico por grupo de estudantes.
A moda se situou no primeiro intervalo, aquele com valores entre 0 e 10 incluindo os extremos do intervalo.
A impressão de assimetria informada pelos coeficientes de assimetria positivos ( Quadro 3.5 ) foi reforçada no gráfico pela concentração no primeiro intervalo. 3.3.2.4 Comentários sobre as respostas à Questão Discursiva 39 Essa questão conduziu o estudante a um contexto histórico referencial ( Constituição de 1988 ), citando o trecho do Código de Ética do Jornalista sobre o direito fundamental do acesso à informação.
Segue -se a descrição de um ato de abuso de autoridade por parte da polícia em comunidade pobre ( prisão ilegal, omissão de informação aos parentes do acusado e proibição de cobertura jornalística sob alegação de que os equipamentos atrapalhariam a ação policial ).
Perguntou -se ao estudante qual deveria ser a atitude do jornalista nesta situação, com base no Código de Ética da profissão.
Esta pergunta foi mais objetiva, apresentando dificuldade baixa.
Mesmo que o estudante não conhecesse o texto do Código, seu conteúdo poderia ser inferido através do enunciado.
A questão possibilitou perceber o grau de desinformação dos estudantes no que concerne ao dever essencial do jornalista : garantir o acesso à informação, informando com idoneidade e sendo informado com liberdade.
A compreensão dos conceitos de liberdade e responsabilidade surgiu na maior parte das respostas submetida à ação repressiva da polícia, não à Lei, à consciência individual ou à compreensão cidadã.
Visivelmente, prevaleceu uma imaturidade social e política na percepção do lugar profissional do jornalista como representante do cidadão numa sociedade democrática.
Em geral, os estudantes responderam sob ângulo de compreensão passiva, evitando posição de enfrentamento com o poder constituído, apesar da caracterização de abuso de autoridade.
Ao mesmo tempo, em muitas respostas, o discurso vago de defesa dos Direitos expressos no Código de Ética refletindo o texto da Constituição de 88 foi recitado sem convicção, em um claro aproveitamento do texto da pergunta.
Além da dificuldade de expressão vernácula e de uma caligrafia sofrível por falta de uso ( poucos escrevem à mão hoje em dia ), as respostas evidenciaram fragilidade de formação do pensamento crítico : uma grande quantidade de respostas justificou o condicionamento da presença da reportagem à permissão das autoridades policiais, em função da necessidade de que a ação local não fosse prejudicada pela presença dos equipamentos.
Isso, apesar do enunciado ser claro : “ Imagine uma investida policial em comunidade pobre, em que houve abuso de autoridade ”.
Muitas são as considerações possíveis em relação a este resultado, no campo da análise sociológica, além da falta de preparo acadêmico para uma bem estruturada autonomia profissional.
Algumas respostas que identificaram na ação policial, o abuso já explicitado na pergunta, mostraram uma atitude reativa, mas ainda frágil do ponto de vista da consciência do papel profissional do jornalista numa situação como a descrita.
Um percentual reduzido, mas felizmente significativo, demonstrou ter compreensão do compromisso profissional com a ética e a cidadania intrínsecas ao jornalismo, atendendo aos critérios corretos da chave de correção. 3.3.2.5 Análise da Questão Discursiva 40 de Componente Específico A questão 40, cujas habilidades aferidas encontram -se descritas no Quadro 3.6, que ainda contém outras informações, seguido da Tabela 3.15 e do Gráfico 3.13, foi aquela em que os grupos de estudantes observados obtiveram o pior desempenho entre as questões discursivas do Componente Específico.
A nota média do grupo dos ingressantes foi 21,57, e a dos concluintes 27,05.
Nesta questão, ocorre também o maior percentual de notas nulas : 54,0 \% entre os ingressantes e 49,7 \% entre os concluintes.
Aqui, a maioria das provas com nota zero decorreu de provas com a questão mal formulada : 58,6 \% entre os ingressantes e 61,1 \% entre os concluintes ( ver Tabela 3.15 ).
Uma pequena minoria ( 1,8 \% ), dentre aqueles com nota nula, utilizou a questão para protestar.
Em ambos os grupos a nota máxima de 100,0 pontos foi alcançada.
Os coeficientes de assimetria positivos indicaram uma inclinação para a esquerda.
O Gráfico 3.13 apresenta a distribuição das notas na questão 40 discursiva do Componente Específico por grupo de estudantes.
A moda se situou no primeiro intervalo, aquele com valores entre 0 e 10 incluindo os extremos do intervalo.
A impressão de assimetria informada pelos coeficientes de assimetria positivos ( Quadro 3.6 ) foi reforçada no gráfico pela concentração no primeiro intervalo. 3.3.2.6 Comentários sobre as respostas à Questão Discursiva 40 Esta questão é mais complexa, já que possibilitou respostas livres, desde que embasadas em estruturas adequadas e lógicas ao enunciado proposto e também a um bom projeto editorial jornalístico.
Verificou -se grande dificuldade dos estudantes em apresentarem um número suficiente de pautas que se adequassem ao número de páginas proposto no enunciado da questão e a indicação de pautas com temas que despertassem realmente o interesse da classe social focada também no enunciado.
Foram consideradas as melhores respostas, as que mantiveram o foco no atendimento ao tema proposto : chances das equipes em 2014 e 2016.
A chave diz que deve haver pertinência da percepção de pauta, tamanhos e formatos, argumentação coerente e lógica para a proposta do caderno impresso.
Foi considerada, portanto para a correção, pontuação para quem adequou as pautas a assuntos concernentes ao tema proposto, qual seja a formação de jovens que virão a participar das equipes.
Alguns exemplos : condições de treino, alimentação, preparação psicológica, patrocínio.
Exemplos de atletas conhecidos que superaram dificuldades materiais no início da área, já que o público era “ classe C ” ( sem indicação para decodificação do estudante, repórter do que significa econômica, social ou culturalmente ) e raras idéias originais e adequadas.
A carência vocabular ficou evidente, assim como a falta de cuidado com a caligrafia.

% matched lemmas: elevado, enfatizar
\end{textsample}
